\documentclass[10pt,a4paper]{amsart}
\usepackage[margin=19mm]{geometry}
\usepackage{amsmath,amssymb,mathtools}
\usepackage[scr=rsfs,cal=euler]{mathalfa}
\usepackage{xfrac}
\usepackage[unicode,hidelinks,bookmarks=false]{hyperref}
\newcommand{\ie}{i.e.\ }
\newcommand{\cf}{cf.\ }
\newcommand{\resp}{respectively\ }
\newcommand{\Subsection}{\subsection}
\providecommand{\nobreakdash}{\nobreak\hbox{-}}
\setlength{\emergencystretch}{2em}
\multlinegap0pt
\usepackage{bm}
\usepackage{bbm}
\usepackage{dsfont}
\usepackage{xspace}
\usepackage{cancel}
\usepackage{mathtools}
\usepackage{cleveref}
\usepackage[shortlabels]{enumitem}
\usepackage{amsthm}
\makeatletter
\@ifundefined{theorem}{\newtheorem{theorem}{Theorem}}{}
\@ifundefined{corollary}{\newtheorem{corollary}[theorem]{Corollary}}{}
\@ifundefined{lemma}{\newtheorem{lemma}[theorem]{Lemma}}{}
\makeatother
\newcommand*{\Naturals}{\mathbb{N}}
\newcommand*{\Reals}{\mathbb{R}}
\newcommand*{\cball}[2]{B_{#2}(#1)}
\newcommand*{\crcdfm}[3]{#1(\cball{#2}{#3})}
\newcommand*{\rd}{\mathrm{d}}
\newcommand{\absval}[1]{\lvert #1 \rvert}
\newcommand{\norm}[1]{\lVert #1 \rVert}
\renewcommand*{\geq}{\geqslant}
\renewcommand*{\leq}{\leqslant}
\title[Strong maximum a posteriori estimation in Banach spaces with]{Strong maximum a posteriori estimation in Banach spaces with Gaussian priors}
\author{Hefin Lambley}
\date{Source: arXiv:2304.13622v3 (CC BY 4.0)}
\begin{document}
\maketitle

\noindent\emph{Source and licence.} Strong maximum a posteriori estimation in Banach spaces with Gaussian priors. Version arXiv:2304.13622v3 (CC BY 4.0), distributed under the Creative Commons Attribution 4.0 International licence, as recorded in the arXiv licence metadata for this exact version and shown on the abstract page https://arxiv.org/abs/2304.13622v3. Full rights notice: licenses/arxiv-2304.13622-CC-BY-4.0.txt.\par\smallskip

\noindent\textit{Editorial scope.} Counted unit: the Lemma whose source label is \texttt{lem:AMF\_limit\_point\_is\_strong} in the article (source0.tex lines 851--854, source label \texttt{lem:AMF\_limit\_point\_is\_strong}), delivered together with the article's own complete original proof of it (source0.tex lines 856--908, one \texttt{proof} environment of 53 lines). No mathematics is written, repaired or re-derived: statement and proof are the article's own text line for line. Required context retained uncounted: eq:Bayesian\_posterior (lines 171--174); thm:main (lines 198--210); cor:Gaussian\_measure\_limits (lines 599--607); lem:AMF\_has\_limit\_point (lines 797--804). Every definition, display and label of those lines is retained unchanged. Omitted from this package and not redistributed: 1--170, 175--197, 211--598, 608--796, 805--850, 855--855, 909--1048. Nothing inside a retained excerpt is omitted or altered: every retained body is delivered exactly as the article's lines are written, so no omission receipt of any kind accompanies this package. Citations: the retained excerpts carry no citation command at all, so no bibliography entry of the article is transcribed and no reference is redistributed. Reference adaptation: none was needed and none was made; every reference inside the delivered unit resolves inside it, so no unresolved reference or citation is produced. Printed slips kept literal: no printed word, symbol, hypothesis or number of the counted statement or of its proof was altered. Layout and class: the article's own class is \texttt{article} with packages including amsfonts, amssymb, amsmath, amsthm, amstext, mathtools, subcaption, float, color, xcolor, graphicx, enumitem, mathrsfs, bm; the wrapper is the lane's pinned \texttt{amsart} editorial wrapper at 10pt on A4, which restates the theorem-like environments the retained text needs and adds the article's own macro definitions that the retained text needs (\texttt{\textbackslash Naturals}, \texttt{\textbackslash Reals}, \texttt{\textbackslash absval}, \texttt{\textbackslash cball}, \texttt{\textbackslash crcdfm}, \texttt{\textbackslash geq}, \texttt{\textbackslash leq}, \texttt{\textbackslash norm}, \texttt{\textbackslash rd}), with extra wrapper packages (bm, bbm, dsfont, xspace, cancel, mathtools, cleveref, enumitem). The delivered numbering is the unit's own. Assets: no retained span contains a figure environment, a graphics-inclusion command, a PDF-page inclusion, a table or a TikZ picture; no image file is copied into this package, the package declares no asset dependency and no image file is redistributed with it. Licence: version arXiv:2304.13622v3 of this article is distributed under the Creative Commons Attribution 4.0 International licence, as recorded in the arXiv licence metadata for this exact version and shown on the abstract page https://arxiv.org/abs/2304.13622v3. A full rights notice is delivered at licenses/arxiv-2304.13622-CC-BY-4.0.txt. Characters: every retained excerpt is pure ASCII, so the delivered engine, XeLaTeX with its default Latin Modern text font, reports no missing character.
\begin{equation}
	\label{eq:Bayesian_posterior}
	\mu^y(\rd x) = \exp\bigl(-\Phi(x)\bigr) \,\mu_0(\rd x).
\end{equation}

\begin{theorem}
	\label{thm:main}
	Let $X$ be a separable Banach space equipped with a centred nondegenerate Gaussian prior $\mu_0$.
	Let $\mu^y$ be the corresponding Bayesian posterior of the form \eqref{eq:Bayesian_posterior} for some continuous potential $\Phi \colon X \to \Reals$, and suppose that, for each $\eta > 0$, there exists $K(\eta) \in \Reals$ such that
	\begin{equation} \label{eq:Phi_lower_bound}
		\Phi(x) \geq K(\eta) - \eta\|x\|_X^{2} \text{~~for all $x \in X$.}
	\end{equation}
	Then:
	\begin{enumerate}[label=(\alph*)]
		\item $\mu^y$ has a strong mode, i.e.\ a strong MAP estimator, and any strong mode lies in the Cameron--Martin space of $\mu_0$;
		\item strong modes, generalised strong modes, weak modes and minimisers of an OM functional for $\mu^y$ coincide.
	\end{enumerate}
\end{theorem}

\begin{corollary}
	\label{cor:Gaussian_measure_limits}
	Let $X$ be a separable Banach space equipped with a centred nondegenerate Gaussian measure $\gamma$.
	Suppose that $(x_{n}, r_{n})_{n \in \Naturals} \subset X \times [0, \infty)$ converges to $(x^{\star}, 0)$.
	After passing to a subsequence without relabelling,
	\begin{equation*}
		\limsup_{n \to \infty} \frac{\crcdfm{\gamma}{x_{n}}{r_{n}}}{\crcdfm{\gamma}{x^{\star}}{r_{n}}} \leq 1.
	\end{equation*}
\end{corollary}

\begin{lemma}[Limit points of AMFs for Bayesian posteriors]
	\label{lem:AMF_has_limit_point}
	Under the assumptions of \cref{thm:main}, if $(x_r)_{r > 0} \subset X$ is an AMF for $\mu^y$, then:
	\begin{enumerate}[label=(\alph*)]
		\item any limit point of $(x_r)_{r > 0}$ lies in $E$;
		\item the net $(x_r)_{r > 0}$ has at least one limit point.
	\end{enumerate}
\end{lemma}

\begin{lemma}[Limit points of AMFs are strong modes]
	\label{lem:AMF_limit_point_is_strong}
	Under the assumptions of \cref{thm:main}, any $X$-strong limit point $x^{\star} \in E$ of an AMF $(x_r)_{r > 0} \subset X$ for $\mu^y$ is a strong mode.
\end{lemma}

\begin{proof}
	Let $x^{\star} \in E$ be some limit point of $(x_{r})_{r > 0}$. 
	To show that $x^{\star}$ is a strong mode, it suffices to check that for any $(r_{n})_{n \in \Naturals} \to 0$,
	\begin{equation} \label{eq:strong_mode_along_subsequences}
		\lim_{n \to \infty} \frac{\mu^{y}(\cball{x^{\star}}{r_{n}})}{M_{r_{n}}} \geq 1.
	\end{equation}
	Indeed, it would be enough to show that any $(r_{n})_{n \in \Naturals} \to 0$ has a further subsequence such that \eqref{eq:strong_mode_along_subsequences} holds along that subsequence:
	this follows from the fact that if $(u_n)_{n \in \Naturals}$ is an arbitrary real sequence and any subsequence of $(u_n)_{n \in \Naturals}$ has a further subsequence converging to $u$, then $u_n \to u$.

	Hence, take any sequence $(r_{n})_{n \in \Naturals} \to 0$; by \cref{lem:AMF_has_limit_point} we may pass to a subsequence of $(r_{n})_{n \in \Naturals}$, which will not be relabelled, such that $(x_{r_n})_{n \in \Naturals}$ converges to some $y^{\star} \in E$.
	As $\Phi$ is continuous, for any $\varepsilon > 0$ there exists $\delta > 0$ such that for any $x \in \cball{y^\star}{\delta}$, it follows that $\absval{\Phi(y^{\star}) - \Phi(x)} < \varepsilon$.
	Since $(r_n)_{n \in \Naturals} \to 0$ and $(x_{r_n})_{n \in \Naturals} \to y^{\star}$, there exists $N \in \Naturals$ such that $\absval{r_n} < \tfrac{\delta}{2}$ and $\norm{x_{r_n} - y^{\star}}_X < \tfrac{\delta}{2}$ for $n \geq N$.
	Hence for such $n$ and any $x \in \cball{x_{r_n}}{r_n}$, we have $\Phi(x_{r_n}) - \Phi(x) < 2\varepsilon$.
	This implies that
	\begin{align}
		\notag
		\frac{\mu^y(\cball{x_{r_n}}{r_n})}{\mu^y(\cball{y^{\star}}{r_n})} &= \exp\bigl(\Phi(y^{\star}) - \Phi(x_{r_n})\bigr) \frac{\int_{\cball{x_{r_n}}{r_n}} \exp\bigl(\Phi(x_{r_n}) - \Phi(x)\bigr)\,\mu_0(\rd x)}{\int_{\cball{y^{\star}}{r_n}} \exp\bigl(\Phi(y^{\star}) - \Phi(x)\bigr)\,\mu_0(\rd x)} \\
		\label{eq:posterior_mass_ratio}
		&< \exp(4\varepsilon) \frac{\mu_0(\cball{x_{r_n}}{r_n})}{\mu_0(\cball{y^{\star}}{r_n})}.
	\end{align}
	By \cref{cor:Gaussian_measure_limits}, we may pass to a further subsequence of $(x_{r_n})_{n \in \Naturals}$ without relabelling such that
	\begin{equation*}
		\limsup_{n \to \infty} \frac{\mu_0(\cball{x_{r_n}}{r_n})}{\mu_0(\cball{y^{\star}}{r_n})} \leq 1,
	\end{equation*}
	and thus taking the $\limsup$ in \eqref{eq:posterior_mass_ratio} and infimising over $\varepsilon > 0$ yields
	\begin{equation*}
		\left(\liminf_{n \to \infty} \frac{\mu^y(\cball{y^{\star}}{r_n})}{\mu^y(\cball{x_{r_n}}{r_n})}\right)^{-1} = \limsup_{n \to \infty} \frac{\mu^y(\cball{x_{r_{n}}}{r_{n}})}{\mu^y(\cball{y^{\star}}{r_{n}})} \leq \limsup_{n \to \infty} \frac{\mu_0(\cball{x_{r_{n}}}{r_{n}})}{\mu_0(\cball{y^{\star}}{r_{n}})} \leq 1.
	\end{equation*}
	Since $(x_{r_{n}})_{n \in \Naturals}$ is a subsequence of an AMF, the previous equation implies that
	\begin{equation*}
		1 \leq \lim_{n \to \infty} \frac{\mu^y(\cball{x_{r_{n}}}{r_{n}})}{M_{r_{n}}} \liminf_{n \to \infty} \frac{\mu^y(\cball{y^{\star}}{r_{n}})}{\mu^y(\cball{x_{r_{n}}}{r_{n}})} = \lim_{n \to \infty} \frac{\mu^y(\cball{y^{\star}}{r_{n}})}{M_{r_{n}}}  \leq 1.
	\end{equation*}
	In particular, the point $x^{\star}$ fixed at the start of the proof is a limit point of the AMF $(x_r)_{r > 0}$, i.e.\ there exists $(s_n)_{n \in \Naturals} \to 0$ such that $(x_{s_n})_{n \in \Naturals} \to x^{\star}$, so the above argument implies the existence of a subsequence such that
	\begin{equation} \label{eq:x_star_maximises_along_sn}
		\lim_{n \to \infty} \frac{\mu^y(\cball{x^{\star}}{s_n})}{M_{s_n}} = 1.
	\end{equation}
	This does not yet hold for \emph{every} sequence $(r_n)_{n \in \Naturals} \to 0$, only the specific sequence $(s_n)_{n \in \Naturals}$.
	To complete the proof, fix an arbitrary $(r_n)_{n \in \Naturals} \to 0$ and $y^{\star} \in E$ as above.
	As $\mu^y$ has an OM functional $I^y$ defined on $E$,
	\begin{equation} \label{eq:OM_functional_rn_sn}
		\lim_{n \to \infty} \frac{\mu^{y}(\cball{x^{\star}}{r_n})}{\mu^{y}(\cball{y^{\star}}{r_n})} = 
		\lim_{n \to \infty} \frac{\mu^{y}(\cball{x^{\star}}{s_n})}{\mu^{y}(\cball{y^{\star}}{s_n})} = 
		\lim_{r \to 0} \frac{\mu^y(\cball{x^{\star}}{r})}{\mu^y(\cball{y^{\star}}{r})} = I^y(y^{\star}) - I^y(x^{\star}).
	\end{equation}
	Hence, using \eqref{eq:x_star_maximises_along_sn} and \eqref{eq:OM_functional_rn_sn}, it follows that
	\begin{align*}
		\lim_{n \to \infty} \frac{\mu^{y}(\cball{x^{\star}}{r_n})}{M_{r_n}} 
		&= \lim_{n \to \infty} \frac{\mu^{y}(\cball{y^{\star}}{r_n})}{M_{r_n}} \lim_{n \to \infty} \frac{\mu^{y}(\cball{x^{\star}}{r_n})}{\mu^{y}(\cball{y^{\star}}{r_n})} \\
		&= \lim_{n \to \infty} \frac{\mu^{y}(\cball{x^{\star}}{s_n})}{\mu^{y}(\cball{y^{\star}}{s_n})}
		\geq \lim_{n \to \infty} \frac{\mu^{y}(\cball{x^{\star}}{s_n})}{M_{s_n}} = 1.
	\end{align*}
	Hence, for any $(r_n)_{n \in \Naturals} \to 0$, there is a further subsequence for which \eqref{eq:strong_mode_along_subsequences} holds, and thus $x^{\star}$ is a strong mode.
\end{proof}

\end{document}
